\scanpage{288}
\mspage{- 102 -}\hfill IX
% archive sheet number: 289
% mounting strip and printed label: 420

Or par le théorème d'orthogonalité appliqué a la fois pour $A_K$ et $A'_K$ on a
\[
V^{\handwritten{\perp}}=W\quad ,\quad (U/V)^{\handwritten{\vee}}\simeq W'(-1)\quad ,
\]
où $W'$ est la partie torique de $U'=T_{\handwritten{\ell}}(A'(\overline K))$, et où le $(-1)$ désigne le twist de Tate. D'autre part, l'homomorphisme (9.2.4) se factorise par $I(\handwritten{\ell})\simeq\mathbb Z_{\handwritten{\ell}}(1)$, de sorte que l'homomorphisme (9.2.4) peut s'interprèter comme un homomorphisme canonique
\[
u:\mathbb Z_{\handwritten{\ell}}(1)\longrightarrow W'\otimes W(-1)\quad ,\tag*{(9.2.5)}
\]
ou encore comme un élément
\[
u\handwritten{\in}W'\otimes W\,(-2)\quad .\tag*{(9.2.6)}
\]

Se rappelant de la définition de $W$ et de $W'$, comme les groupes de points à valeurs dans $\overline K$ des deux faisceaux $\handwritten{\ell}$-adiques constants tordus $T_{\handwritten{\ell}}(A^{\unsure{o|0}})^{t}$ et $T_{\handwritten{\ell}}(A'^{o})^{t}$, dont les fibres spéciales sont respectivement $T_{\handwritten{\ell}}(T_o)$ et $T_{\handwritten{\ell}}(T'_o)$, on trouve les isomorphismes canoniques
\[
\left\{\begin{aligned}
W&\simeq T_{\handwritten{\ell}}(T_o)(\overline k))\simeq\overset{\handwritten{\scriptscriptstyle\vee}}{M}\otimes\mathbb Z_{\handwritten{\ell}}(1)\\
W'&\simeq T_{\handwritten{\ell}}(T_o(\overline k))\simeq M'\otimes\mathbb Z_{\handwritten{\ell}}(1)\quad ,
\end{aligned}\right.\tag*{(9.2.7)}
\]
$M$ et $M'$ désignant les groupes des caractères relativement à $\overline k$ de $T_o$ et de $T'_o$ respectivement :
\[
T_o=D(M)\quad ,\quad T'_o=D(M')\quad .\tag*{(9.2.8)}
\]
Par suite, l'élément $u$ de (9.2.6), qui exprime l'opération du groupe d'inertie $I$ sur le module de Tate $U$ de $A_K$, s'interprète comme un élément
